\documentstyle[%


righttag%


,11pt%


,amssymb%


,russian%               remove in Eng.


,izvru%                 Set izven% in Eng.


% ,izvru-ks             for brief communications


]{amsart}





% l_o ?





%\newcommand{\ov}{\overline}


%\newcommand{\wt}{\widetilde}


\newcommand{\Spin}{\operatorname{Spin}}


% NB: Use! \marginpar{\tiny{line 60}}


\newcommand{\tts}[1]{}


\renewcommand{\baselinestretch}{1.06}





%\hoffset=-15mm%                  For Eng.





\setcounter{page}{9}


\god{1998}


\nomer{6~(433)} %                        Remove in Eng.


%\nomer{Vol.\,42, No.\,6}%               Set in Eng.


%\str{\pageref{begin}--\pageref{end}}%  Set in Eng.


\udk{530.12}


\author{\it Р.Ф.\,Билялов, Б.С.\,Никитин}


% NB:   ^^ remove in Eng.


\title{Спиноры в произвольных реперах.\\


Ковариантная производная \\


и производная Ли спиноров}





\date{}


%\pagestyle{myheadings}%         Set in Eng.





\pagestyle{plain} %              Remove in Eng.


\begin{document}





%\thispagestyle{plain}%          Set in Eng.





\maketitle


\label{begin}





\section*{Введение}





В наиболее общей математической форме спиноры были введены Э.\,Картаном


в 1913~г. [1], а  затем в 1928~г. переоткрыты Ван дер Варденом [2]


в связи с физическими исследованиями Дирака. Спиноры дают линейное


представление группы вращений пространства n измерений, причем каждый


спинор определен при помощи $2^\nu$ составляющих ($n=2\nu+1$ или


$2\nu$). Спиноры четырехмерного пространства входят в знаменитое


уравнение Дирака для электрона, причем четыре волновых функции


являются ни чем иным, как составляющими спинора.





Теория спиноров в римановом пространстве-времени общей теории


относительности была развита в работах [3]--[5].


Симметрический тензор


энергии-импульса для спинорных полей с произвольными функциями


Лагранжа был получен Л.\,Розенфельдом в 1940~г. [6]. В работе


Розенфельда рассуждения при применении теоремы Н\"eтер к построению


законов сохранения не носят ковариантный характер, ковариантная


запись законов сохранения достигается с помощью остроумных


преобразований. Ковариантное построение симметрического тензора


энергии-импульса для спинорных полей получено в [7] путем применения


производной Ли спиноров, предложенной И.\,Косман [8] и обоснованной с


теоретико-групповой точки зрения в [9]. Теоретико-групповое


обоснование производной Ли спиноров требует рассмотрения спиноров в


произвольных неортогональных реперах. Поэтому возникает проблема


построения спинорного анализа в произвольных реперах, в первую


очередь проблема построения ковариантной производной и производной Ли


в произвольных реперах.





\section{Группа Spin(4) [10]}





Рассмотрим пространство-время Минковского, отнесенное к некоторой


ортогональной декартовой системе координат, и обозначим через


$\eta_{\alpha\beta}$


($\eta_{\alpha\beta}=0$, $\alpha\ne\beta$; $\eta_{00}=1$,


$\eta_{11}=\eta_{22}=\eta_{33}=-1$) его метрический тензор. Здесь и далее


все индексы, кроме


заглавных латинских, принимают значения 0,~1,~2,~3; заглавные


латинские --- 1,~2,~3,~4. Введем систему четырех $4\times4$ матриц


Дирака $\gamma_\alpha$ с комплексными элементами


$\gamma_\alpha=(\gamma_\alpha{}^A{}_B)$,


$(\gamma^\alpha=\eta^{\alpha\beta}\gamma_\beta)$,


удовлетворяющих соотношению


$$


\gamma_\alpha\gamma_\beta+\gamma_\beta\gamma_\alpha=


2\eta_{\alpha\beta}I,


$$


где $I$ --- единичная $4\times4$ матрица. Матрицы Дирака определяются


с точностью до выбора базиса в 4-мерном комплексном пространстве


и в теоретической физике их обычно выбирают в виде (представление


Дирака)


\begin{alignat}{2}


\gamma_0&=


\begin{pmatrix}


1 & 0 & 0 & 0 \\


0 & 1 & 0 & 0 \\


0 & 0 & -1 & 0 \\


0 & 0 & 0 & -1


\end{pmatrix}


, &\qquad


\gamma_1&=


\begin{pmatrix}


0 & 0 & 0 & -1 \\


0 & 0 & -1 & 0 \\


0 & 1 & 0 & 0 \\


1 & 0 & 0 & 0


\end{pmatrix}


,\notag \\


\gamma_2&=


\begin{pmatrix}


0 & 0 & 0 & i \\


0 & 0 & -i & 0 \\


0 & -i & 0 & 0 \\


i & 0 & 0 & 0


\end{pmatrix}


, &\qquad


\gamma_3&=


\begin{pmatrix}


0 & 0 & -1 & 0 \\


0 & 0 & 0 & 1 \\


1 & 0 & 0 & 0 \\


0 & -1 & 0 & 0


\end{pmatrix}


.


\end{alignat}





Матрицы $I$, $\gamma_\alpha$, $\gamma_\alpha\gamma_\beta$,


$\gamma_\alpha\gamma_\beta\gamma_\gamma$,


$\gamma_\alpha\gamma_\beta\gamma_\gamma\gamma_\delta$


($\alpha \ne\beta \ne\gamma\ne\delta$)


линейно независимы; в их линейной оболочке над полем вещественных


чисел совокупность неособенных матриц образует группу Ли $\Gamma$


относительно обычного умножения матриц. Подмножество


$G\subset\Gamma$,


порожденное матрицами $\Lambda$ с $\det\Lambda=1$, для


которых линейная оболочка $M$ матриц $\gamma_\alpha$ с вещественными


коэффициентами остается инвариантной при преобразовании


$\Lambda{M}\Lambda^{-1}=M,$ тоже представляет группу Ли, она называется


группой $\Spin(4)$. Если положить


$\Lambda\gamma_\alpha\Lambda^{-1}=A^{\beta}_{\alpha}\gamma_{\beta}$,


то матрица $A=(A^{\beta}_{\alpha})$ представляет преобразование


Лоренца в $L(4)$.


Группы $ L(4)$ и $\Spin(4)$ локально изоморфны, группа $\Spin(4)$


гомоморфна группе $L(4)$, ядро гомоморфизма состоит из матриц


$\pm{I}$, в итоге группа $\Spin(4)$ дважды накрывает группу $L(4)$.


Если преобразование Лоренца $L$ из связной компоненты единицы


имеет вид $L=\exp(K)$, $K-\eta$ --- кососимметрическая матрица:


$\eta K=-K^{T}\eta$, то соответствующее преобразование $\Lambda$


из $\Spin(4)$ имеет вид $\Lambda=\pm\exp(\gamma_{\alpha}


\gamma^{\beta}K^{\alpha}_{\beta})$.


При выборе матриц Дирака в виде (1) для временного и


пространственного отражений в пространстве Минковского соответствующие


преобразования соответственно имеют вид


$\eta(T)\gamma^{1}\gamma^{2}\gamma^{3}$ и $\eta(P)\gamma^{0}$,


где $\eta^2(T)=\pm1$, $\eta^2(P)=\pm1$.





\section{Римановы и спинорные многообразия}





Пусть $V$ есть пространственно-временное риманово многообразие с


метрикой


$ds^2=\linebreak g_{\alpha\beta}(x)dx^{\alpha}dx^{\beta}$


сигнатуры $(+,-,-,-)$ в локальных координатах $x^\alpha$. Через


$E(V)$


обозначим главное расслоение линейных реперов многообразия $V$, а


через $O(V)$ --- главное расслоение ортонормированных реперов с группой


$L(4)$ в качестве структурной группы. Предполагаем, что $V$


ориентируемо, т.\,е. $O(V)$ допускает сужение структурной группы


$L(4)$ до $L_{0}(4)$-связной компоненты единицы в $L(4)$. Репер


в точке $x$ на $ V$


будем обозначать $e(x)$, он задается четырьмя линейно независимыми


векторными полями, имеющими в локальных координатах вид


$e_a{}^\alpha(x)$, где $\alpha$ --- координатный индекс, $a$ --- номер


вектора репера. Наряду с репером $(e_a{}^\alpha)$ рассматриваем и


корепер $e^a{}_\alpha(x)$, которому соответствует базис


$\Theta^a=e^a{}_\alpha(x){dx}^\alpha$ в пространстве 1-форм. Если


репер $e=(e_a{}^\alpha )$ ортонормированный, то


$e_a{}^\alpha e_b{}^\beta g_{\alpha \beta}(x)=\eta_{ab}$,


где $\eta_{ab}=(1,-1,-1,-1)$.


Локальные координаты в $E(V)$ можно построить следующим образом.


Репер $(e_a{}^\alpha)$ называется натуральным,


если линейные операторы


$e_a{}^\alpha\frac{\partial}{\partial x^\alpha}$


имеют вид


$e_a{}^\alpha\frac{\partial}{\partial x^\alpha}=


\frac{\partial}{\partial x^a}$,


т.\,е.


$e_a{}^\alpha=\delta^\alpha_a$. Для произвольного репера $e(x)$ получим


$e(x)_a{}^\alpha=\delta^\alpha_b A^b_a$,


и локальные координаты репера $e(x)$ будут иметь вид


$(x,A)$, где $A=(A^b_a)$. Для натурального репера будем использовать


обозначение $\partial x$. Если


$\partial y$ --- натуральный репер для карты с координатами


$y^\alpha$, то


$\partial y=\frac{\partial y}{\partial x}\partial x$,


где


$\frac{\partial y}{\partial x}=\Bigl( \frac{\partial y^\alpha}


{\partial x^\beta}\Bigr)$ --- якобиева матрица.


Если в новой карте ``отсчет'' реперов начинать снова с натурального,


то в области пересечения двух карт с локальными координатами $x$


и $y$ в координатах $y$ имеем $e(x)\rightarrow(y(x),A(y(x)))$,


т.\,е. функции перехода $\psi_{yx}$ представляют собой постоянные


единичные $4\times4$ матрицы $I$. В случае $O(V)$ тоже можно


подобрать атлас, для которого функции перехода равны единичной


матрице, если в $O(V)$ существует сечение $l_o:V\rightarrow{O(V})$,


$x\rightarrow{l_o(x)}$ такое, что на пересечении карт имеем


$l_o(y)\rightarrow{\frac{\partial y}{\partial x}l_o(x)}$. К


каждому ортонормированному реперу отнесем также матрицы Дирака,


внешние индексы которых уже будем обозначать строчными буквами


латинского алфавита.





Наряду с главным расслоением $O(V)$ ортонормированных реперов


можно ввести главное расслоение $S(V)$ спинорных реперов со структурной


группой $\Spin(4)$, которое можно получить из $O(V)$ путем


расширения $O(V)$. Поэтому расслоение спиноров можно считать


ассоциированным с расслоением $O(V)$ ортогональных реперов или с


расслоением $S(V)$ спинорных реперов.





\section{Ковариантная производная спиноров}





Ковариантная производная от спинора впервые получена В.А.\,Фоком в


[3]. Для получения ковариантной производной спинора использовался


тот факт, что билинейные комбинации спиноров


$A_i=\ov\psi{\alpha_i}{\psi}$, где $\alpha_0=I$,


${\alpha_i={\gamma^5}{\gamma^i}{\gamma^0}}$ $(i=1,2,3)$,


$\alpha_4={-i\gamma^2}$, ${{\alpha_5}={\gamma^0}{\gamma^2}}$,


образуют один


четырехвектор $(i=0,1,2,3)$ и два инварианта ${i=4,5}$,


для которых закон параллельного переноса при бесконечно малых


преобразованиях известен. Предполагая,


что приращение спинора при таком преобразовании пропорционально


исходному спинору с некоторой матрицей пропорциональности и


дифференциалу смещения, т.\,е.


${\delta\psi=\sum{e_lC_l{ds}_l\psi}}$, где $C_l$ --- матрица


пропорциональности, ${ds}_l$ --- ортогональные компоненты смещения


(и неявно предполагая выполнения правила Лейбница),


сравнивая законы преобразования $A_i$ с законом преобразования вектора


и инварианта, получим систему матричных уравнений на матрицу


пропорциональности.


${C^{+}{_l}{\alpha_i}+{\alpha_i}{C_l}={\sum{e_k}\gamma_{ikl}


\alpha_k}}$


($i=0$--3) и ${C^{+}{_l}{\alpha_i}+{\alpha_i}{C_l}=0}$


(для $i=4,5).$


Общее решение этих уравнений содержит наряду с хорошо известной


связностью Фока аддитивный член, пропорциональный единичной матрице,


который интерпретируется Фоком как потенциал электромагнитного поля.


Этот произвол в определении ковариантной производной спинора


устраняется путем наложения следующего требования: ковариантная производная


метрического спинора в пространстве спиноров должна быть равна


нулю [11].





Изящный способ построения ковариантной производной спиноров предложен


А.\,Лих\-не\-ро\-ви\-чем [10]. На римановом многообразии существует


связность, определяемая символами Кристоффеля от метрического тензора и


называемая римановой связностью. Связность на расслоении линейных


реперов называется линейной связностью, а на расслоении ортогональных


реперов --- лоренцевой. Лоренцева связность тогда и только тогда


является римановой связностью, когда ассоциированная с ней


линейная связность обладает нулевым кручением и 1-форма


$\omega=(\omega^a{}_b)$ при заданном сечении $e_o(x)$ на


главном расслоении $O(V)$ имеет вид


$$


\omega^a{}_b(\xi)=-c^a{}_{bc}\xi^c,


$$


где $c_{abc}$ --- коэффициенты вращения Риччи, $\xi^\alpha$


--- векторное поле.


А.\,Лихнерович этой форме связности на $O(V)$ ставит в соответствие


спинорную связность $\sigma$ на $S(V)$ по закону


$\sigma(V_z)=p'{}^{-1}\omega(pV_z)$, где $p(V_z)$ --- проекция на


$O(V)$ касательного вектора $V_z$ к $z$ на $S(V)$ и $p'$ ---


изоморфизм алгебры Ли группы $\Spin(4)$ на алгебру Ли $L(4)$.


После введения некоторого локального сечения в $S(V)$ форма


связности $\sigma$ оказывается равной


$$


\sigma=\frac{1}{4}\omega^a{}_b\gamma_a\gamma^b.


$$


Если $M$ --- некоторое 4-мерное комплексное векторное пространство,


на котором действует $\Spin(4)$, то контравариантный 1-спинор


$\psi$ в точке $x\in V$ определяется как отображение


$z\rightarrow{\psi(z)}$ из $\pi^{-1}(x)$ в $M$ такое, что


$$


\psi(z\Lambda^{-1})=\Lambda\psi(z) \quad (\Lambda\in \Spin(4)).


$$


Ковариантная производная спинора


$\xi^\alpha\nabla_\alpha\psi$


в направлении вектора $\xi^\alpha$ определяется формулой


\begin{equation}


\xi^\alpha\nabla_\alpha\psi=


\xi^\alpha\partial_\alpha\psi+\frac{1}{4}\omega^a{}_b(\xi)


\gamma_a\gamma^b\psi.


\end{equation}


В дальнейшем, следуя [12], будем рассматривать спиноры как элементы


спинорного расслоения, ассоциированного с главным расслоением


$O(V)$ ортогональных реперов. Ассоциированное расслоение со


стандартным слоем $M$, на котором группа $\Spin(4)$, как


представление группы $L(4)$, действует правосторонним образом


$$


L\in L(4),\quad L:\ M\rightarrow M,\quad \psi\rightarrow \Lambda


^{-1}(L)\psi=\psi\cdot L,


$$


определяется как фактор-пространство тензорного произведения


$O(V)\times M$ относительно действия группы $L(4)$. Пусть для


каждого $x\in V$ представитель класса эквивалентности выбран


в виде $(x,L(x),\psi)$, где $L(x)$ задает некоторое поле


реперов, a $\psi\in M$. Если $\psi$ пробегает $M$, то


${(x,L(x),\psi)}$ будет определять слой $M_x$ в точке $x$.


Каждый элемент $z=(x,L_1)\in O(V)$ порождает отображение


$$


M\rightarrow M_x,\quad \psi\rightarrow z\cdot \psi=


(x,L(x),\psi\cdot (L_1{}^{-1}L(x))).


$$


Кривая


$$


t \rightarrow (x(t),L(x(t))(1-t\omega (\xi)) \cdot \psi(x)=


(x,L(x(t),\psi\cdot(1+t\omega(\xi)))=


\biggl(x,L(x(t)),\biggl(1-\frac{t}{4}\omega^a{}_b(\xi)\gamma_a\gamma^b


\biggr)\biggr)\psi


$$


представляет собой горизонтальный путь в ассоциированном расслоении


и определяет параллельный перенос $\psi(x)$ из точки $x$ в


точку $x(t)$. Поэтому ковариантная производная


$\nabla_\xi\psi$


в направлении вектора $\xi$ есть


$$


\nabla_\xi\psi=\lim_{t \to 0}\frac


{\psi(x(t))-\psi(x) \cdot (1+t\omega(\xi))}{t}=


\xi^\alpha\partial_\alpha\psi+\frac{1}{4}\omega^a{}_b(\xi)


\gamma_a\gamma^b\psi.


$$


Мы приходим к тому же результату (2).





\section{Производная Ли спинора}





Производную Ли спинорных полей будем строить по аналогии с


производными Ли векторных полей ([12], с.\,36). С помощью бесконечно


малого преобразования координат $x\rightarrow x'=x+t\xi$


векторному полю $A^\alpha (x-t\xi)$


в точке $x-t\xi$ можно сопоставить увлеченное векторное поле


в точке $x$ по правилу $\wt A^\alpha(x)=\frac{\partial


x'{}^\alpha}{\partial x^\beta}A^\beta (x-t\xi).$


Производная Ли векторного поля $A^\alpha$ вдоль векторного поля


$X=\xi^\alpha\partial_\alpha$ определяется формулой


$$


L_X A^\alpha=\lim_{t \to 0}(A^\alpha(x)-\wt A^\alpha(x))/t=


[X,A]^\alpha,


$$


где $A = A^\alpha \partial_\alpha$.


При эквивалентном определении производной Ли векторного поля


на языке реперов реперу $e(x-t\xi)$ в точке $x-t\xi$ нужно


сопоставить увлеченный репер в точке $x$


$$


\wt e (x)=(\wt e_a{}^\alpha),\quad \wt e_a{}^\alpha=


e_a{}^\beta(x-t\xi)\frac{\partial x'{}^\alpha}


{\partial x^\beta}=


(\delta^b_a +tK_a{}^b+\cdot)e_b{}^\alpha(x),\quad K_a{}^b=


-e^b{}_\alpha L_X e_a{}^\alpha.


$$


В матричных обозначениях $\wt e=e(1+tK+\cdot)$. Нужно считать,


что при переходе


от $e(x-t\xi)$ к $\wt e(x)$ реперные компоненты не


изменяются, увлеченные значения $\wt A^a(x)$ получаются с


помощью преобразования $(1+tK+\cdot)$


при переходе от $\wt e(x)$ к $e(x)$


$$


\wt A^a(x)=A^b(x-t\xi)(\delta^a_b+tK_b{}^a+\cdot).


$$


Если положить


$$


L_xA^a=\lim_{t\to0}{(A^a(x)-\widetilde A^a(x))}/t,


$$


то получим


\begin{equation}


L_XA^a=\xi^\alpha\partial_\alpha A^a-K_b{}^a A^b=


X(A^a)-K_b{}^aA^b=e^a{}_\alpha L_XA^\alpha.


\end{equation}


Следует заметить, что при таком подходе к определению производных Ли


вектора производные Ли векторов репера оказываются равными


нулю. Векторы репера $e_a{}^\alpha$ характеризуются двумя индексами


различной природы, каждому из которых нужно сопоставлять


соответствующие операции производной Ли, а их суммарное действие


приводит в итоге к нулевому значению. Если $e(x)$ --- ортрепер,


то соответствующее преобразование $(1+tK+\cdot)$ в общем случае


произвольного вектора $\xi^\alpha$ не является преобразованием Лоренца,


поэтому, чтобы получить производную Ли спинора, И.\,Косман [8]


использует преобразование $K_A=(K-K^*)/2$, где $K^*$ --- $\eta$-сопряженное


преобразование к $K$. Матрица $K$ в (3) для случая спиноров заменяется на


$-\frac{1}{4}K_{Ab}{}^a\gamma_a\gamma^b$.


Тогда для производной Ли спинора получается выражение


\begin{equation}


L_X\psi=X(\psi)-\frac{1}{4}K_{Ab}{}^a\gamma_a\gamma^b\psi.


\end{equation}


При таком определении производной Ли спиноров невозможно


проводить последовательно строгие рассуждения по применению


теоремы Н\"етер к выводу законов сохранения для спинорных полей.


Нужно уметь преобразовывать спиноры при нелоренцевых


преобразованиях репера или, что то же самое, уметь рассматривать


спиноры в произвольных, неортонормированных реперах.





\section{Спиноры в произвольных реперах}





Для того чтобы иметь возможность рассматривать спиноры в


произвольных реперах, необходимо задавать закон преобразования


спиноров при переходе от одного репера к другому с помощью


произвольного линейного преобразования $A$, т.\,е. необходимо


представление группы Лоренца $L(4)$ в пространстве спиноров расширить


до представления полной линейной группы $GL(4)$. Для


этого воспользуемся полярным разложением в пространстве Минковского


[9]: всякое неособенное линейное преобразование $A$ в пространстве


Минковского единственным образом представимо в виде произведения


симметрического преобразования $S$ ${(S^T\eta=\eta S)}$ и


преобразования Лоренца $L$ или в виде произведения симметрических


преобразований $T$ и $S$ и преобразования Лоренца $L$.


Преобразование $S$ имеет то же самое корневое подпространство,


что и преобразование ${\eta^{-1}A^{T}\eta A}$, собственные значения


преобразования $S$ положительны или имеют положительную вещественную


часть. Матрица преобразования $T$ в каноническом базисе преобразования


${\eta^{-1}A^{T}\eta }A$ имеет вид


$$


T=


\begin{pmatrix}


0 & 1 & 0 & 0 \\


1 & 0 & 0 & 0 \\


0 & 0 & 1 & 0 \\


0 & 0 & 0 & 1


\end{pmatrix}


.


$$


Если $e=(e_a{}^\alpha)$ --- произвольный репер в пространстве


Минковского, то $g_{ab}=e_a{}^\alpha e_b{}^\beta\eta_{\alpha\beta}$


--- реперные компоненты метрического тензора. Если репер $e$


ортонормированный, то в этом случае будем применять для $g_{ab}$


обозначение $\eta_{ab}$.





Для каждого репера $e$ в пространстве Минковского, характеризуемого


метрическим тензором, существует единственное преобразование $P$,


которое есть или симметрическое преобразование или произведение


симметрических преобразований $T$ и $S$ и которое удовлетворяет


условию $P^T gP=\eta$.





Для того чтобы расширить представление группы Лоренца $L(4)$ до


представления $GL(4)$, следуя [9], спинор $\psi$ (пусть для


определенности это контравариантный спинор типа $(1,0)$), приписанный


к реперу $e$, будем рассматривать вместе с преобразованием $P$,


которое метрический тензор $g$, приписанный к тому же реперу,


переводит в $\eta:P^T gP=\eta$. Таким образом, спинор $\psi$,


отнесенный к реперу $e$, будет характеризоваться индексом


$P:\psi_P$, что будем обозначать с помощью пары $(\psi,P)$. Если


теперь преобразование $A$ переводит репер $e'=(e'_a{}^\alpha)$


в репер


$e=(e_a{}^\alpha):e=e'A,\; e_a{}^\alpha=e'_b{}^\alpha A_a{}^b$,


то действие преобразования, соответствующего линейному преобразованию


$A$, на пару $(P,\psi)$, обозначаемое как


$A\cdot(P,\psi)$, определено следующим образом:


$$


A\cdot(P,\psi)=(P',\psi'),


$$


где преобразование $P'$ приписано к реперу $e'$,


$\psi'=\Lambda (L)\psi$, $L=P'{}^{-1}AP$, а $\Lambda(L)$ ---


спин-преобразование, соответствующее преобразованию Лоренца $L$.





Мы получили представление полной линейной группы $GL(4)$ в


пространстве пар $(P,\psi)$. Чтобы убедиться в этом, нужно показать,


что при рассмотрении перехода $e'\rightarrow e''\cdot B$


выполняется соотношение


\begin{equation}


B\cdot (A\cdot(P,\psi))=(BA)\cdot (P,\psi).


\end{equation}


Положим


$$


B\cdot (P',\psi')=(P'',\psi''),


$$


где, по определению, $P''$ соответствует реперу $e''$, а


$\psi''=\Lambda(L')\psi'$ при $L'=P''{}^{-1}BP$. Имеем


\begin{gather*}


e=e''(BA),\\


\psi''=\Lambda (L')\psi'=\Lambda(L')\Lambda(L)=\Lambda(L'L)


\psi=


\Lambda(P''{}^{-1}BP'{}^{-1}P'AP)\psi=


\Lambda(P''{}^{-1}(BA)P)\psi,


\end{gather*}


что означает (5). Если ограничиться рассмотрением только


ортонормированных реперов $e$, $e'$, то $P=P'=1$ (1 --- тождественное


преобразование в пространстве Минковского), $A=L$ , т.\,е. $A$


есть преобразование Лоренца, и мы получаем обычное представление


группы Лоренца в пространстве спиноров


$$


A\cdot (1,\psi)=(1,\Lambda(A)\psi).


$$


Закон преобразования сопряженных (ковариантных) спиноров при


преобразовании $e'\rightarrow e=e'A$ будет иметь вид


$$


\ov\psi'=\ov \psi\Lambda^{-1}(L).


$$


Каждому реперу $e$ приписываются и матрицы $\gamma^a$ (их снова


назовем матрицами Дирака), преобразующиеся по закону


$$


\gamma'{}^a=A_b{}^a\Lambda(L)\gamma^b\Lambda^{-1}(L)


$$


и удовлетворяющие перестановочным соотношениям


$$


\gamma^a\gamma^b+\gamma^b\gamma^a=2g^{ab} \cdot1.


$$


Величины $\ov \psi\gamma^a\psi$ при преобразованиях реперов


преобразуются по векторному закону


$$


\ov \psi'\gamma'{}^a\psi'=A_b{}^a\ov \psi\gamma^b\psi.


$$





Но мы к каждому реперу будем приписывать обычные матрицы Дирака


(1), т.\,к. в дальнейшем матрицы Дирака используются


только при рассмотрении таких спин-преобразований, в которых матрицы Дирака


считаются обычными.





\section{Ковариантная производная спиноров в произвольных реперах}





Пусть $E(V)$ --- главное расслоение линейных реперов со связностью


$\omega$, ассоциированной c римановой связностью на $O(V)$ и


соответствующей локальному сечению $z(x)=(x,A(x))$, $x\in V$,


$A(x)\in GL(4)$. Если $x(t)=x+t\xi$ --- путь в $V$, порожденный


векторным полем $\xi^\alpha$, то соответствующий горизонтальный


путь для малых значений $t$ будет представлять собой кривую


$t\rightarrow (x(t),A(x(t))\cdot (1-t\omega(\xi)))$. Присоединенное


расслоение со стандартным слоем $F$, на котором $GL(4)$ действует


правосторонним образом:


$$


A\in GL(4),\quad A: F\rightarrow F, \quad \Psi\rightarrow \Psi\cdot A,


$$


определяется как фактор-пространство тензорного произведения


относительно действия группы $GL(4)$. Пусть для каждого $x\in V$


представитель класса эквивалентности выбран в виде


$(x,A(x),\Psi)$, где $A(x)$ задает некоторое поле


реперов, $\Psi\in F$. Если $\Psi$ пробегает $F$, то


$\{(x,A(x),\Psi)\}$ будет представлять собой слой $F_x$ в точке $x$.


Каждый элемент $z=(x,B)$ порождает отображение


$$


F\rightarrow F_x,\quad \Psi\rightarrow z\cdot\Psi=


(x,A(x),\Psi\cdot(B^{-1}A(x))).


$$


Кривая


$$


t\rightarrow (x(t),A(x(t))(1-t\omega(\xi))).\Psi (x)=


(x(t),A(x(t)),\Psi(x)\cdot(1+t\omega(t)))


$$


является горизонтальным путем в присоединенном расслоении и


определяет параллельный перенос $\Psi(x)$ из точки $x$ в точку


$x(t)$. Поэтому ковариантная производная


$\nabla_\xi\Psi$ в направлении вектора $\xi$


есть


$$


\nabla_\xi\Psi=\lim_{t\to 0}\frac{\Psi(x(t))-


\Psi(x)\cdot(1+t\omega(\xi))}{t}=\xi^\alpha\partial_\alpha\Psi


-\Psi(x)\cdot\omega(\xi).


$$


В нашем случае $\Psi=(P,\psi)$. Левое действие


$A: \Psi\rightarrow A\cdot\Psi$ преобразуется в правое по правилу


$A:\Psi\rightarrow A^{-1}\cdot\Psi=\Psi'$, $\Psi'=(P',\psi')$,


$g'=A^TgA$, $P'{}^Tg'P'=\eta$, $\psi'=\Lambda(L')\psi$,


$L'=P'{}^{-1}A^{-1}P$.


Пусть $P$ --- симметрическое преобразование.


Если предположить, что при преобразовании $e'=e(1+t\omega)$ имеет место


$g'=(1+t\omega)^Tg(1+t\omega)=g+t(\omega^Tg+g\omega)$, т.\,е.


$\delta g=\omega^Tg+g\omega=\partial_\xi g$, $P'=P+t\delta P$,


где $\delta P$ --- $\eta$-симметрическое преобразование, то


$\delta P$ должно удовлетворять уравнению


$$


(\delta P)^TgP+P^Tg\delta P=-P^T\delta gP.


$$


Если учесть, что $P^TgP=\eta$, то это уравнение запишется в виде


\begin{equation}


P\delta P+\delta PP=\delta Q,


\end{equation}


где


$\delta Q=-P^2\eta^{-1}\delta gP^2$ ---


симметрическое преобразование. Решение этого уравнения


будем искать в виде


$$


\delta P=\sum_{m,n=0}^{\infty}a_{mn}(P-1)^m\delta Q(P-1)^n,


$$


где $a_{mn}=a_{nm}$. Симметричность $a_{mn}$ по индексам $m$ и $n$


обеспечивает $\eta$-симметричность


$\delta P$. Для $a_{mn}$ получаем систему уравнений


\begin{gather*}


2a_{00}=1, \\


2a_{m0}+a_{m-1\ 0}=0,\quad m\ge 1, \\


2a_{mn}+a_{m-1\ n}+a_{m\ n-1}=0,\quad m,n\ge 1,


\end{gather*}


решение которой имеет вид


$$


a_{mn}=\frac{(-1)^{m+n}}{2^{m+n+1}}C^m_{m+n}.


$$


Для $L'=1+t\delta L$, $\psi'=\Lambda (L')=\psi+t\delta \psi$ находим


$$


\delta L=-P^{-1}\delta P- P^{-1}\omega P,\quad


\delta \psi=\frac{1}{4}


\gamma_a\gamma^b(\delta L)^a{}_b\psi.


$$


В итоге для ковариантной производной от $\Psi=(P,\psi)$ получаем


\begin{gather*}


\nabla_\xi P=\xi^\alpha\partial_\alpha P-\delta P=


\xi^\alpha\partial_\alpha P-\sum_{m,n=0}^{\infty}


\frac{(-1)^{m+n}}{2^{m+n+1}}C^m_{m+n}(P-1)^m\delta Q (P-1)^n, \\


\nabla_\xi\psi=\xi^\alpha \partial_\alpha \psi +


\frac{1}{4}\gamma_a\gamma^b(P^{-1}\delta P+P^{-1}\omega P)^a{}_b\psi.


\end{gather*}


В ортонормированном репере, где $P=I$, $g=\eta$ и $\partial_\xi g =0$,


имеем


$\nabla_\xi P=0 $, а $\nabla_\xi\psi $


определяется формулой (2). Оказывается, вообще,


$\nabla_\xi P=0$ ввиду того, что


$\partial_\xi P=\delta P$. Действительно, из $P^T gP=\eta$,


дифференцируя по направлению $\xi^\alpha$, получаем уравнение


$\partial_\xi PP^{-1}+P^{-1}\partial_\xi =


-P\eta^{-1}\partial_\xi gP$, которое эквивалентно уравнению (6),


поэтому $\partial_\xi P=\delta P $.





Докажем, что при любых преобразованиях репера $e=e'A$


ковариантная производная $\nabla_\xi\psi$


тоже есть спинор, т.\,е. $\nabla_\xi\psi'=


\Lambda\nabla_\xi\psi $. Для связности $\omega$


при преобразованиях репера имеем закон преобразования


$$


\omega'=A\omega A^{-1}+A\partial_\xi A^{-1}.


$$


Преобразование Лоренца $L$ и соответствующее спин-преобразование


$\Lambda$ связаны уравнением


$\Lambda\gamma\Lambda^{-1}=L^{-1}\gamma$, где


$\gamma=(\gamma^a).$


Дифференцируя это уравнение в направлении вектора $\xi^\alpha$,


находим


$$


\Lambda^{-1}\partial_\xi\Lambda\gamma-


\gamma \Lambda^{-1}\partial_\xi \Lambda=-L^{-1}\partial_\xi L


\gamma,


$$


откуда


$$


\partial_\xi\Lambda=\frac{\Lambda}{4}\gamma^TL^{-1}


\partial_\xi L\gamma,


$$


где $\gamma^T=(\gamma_a)$. Если воспользоваться тем, что


$A=P'LP^{-1}$, то находим


$$


\delta L'=-P'{}^{-1}\partial_\xi P'-P'{}^{-1}\omega' P'=


\partial_\xi LL^{-1}+L\partial_\xi P^{-1}PL^{-1}-LP^{-1}


\omega PL^{-1}.


$$


Подставляя это выражение в $\nabla_\xi\psi',$ имеем


$$


\nabla_\xi\psi'=


\Lambda \nabla_\xi\psi+\partial_\xi \Lambda \psi+


\frac{\Lambda}{4}\gamma_T \delta L\gamma \psi-


\frac{1}{4}\gamma_T(\partial_\xi LL^{-1}+L\partial_\xi P^{-1}


PL^{-1}-LP^{-1}\omega P)\gamma \Lambda \psi.


$$


В последнем слагаемом правой части, пользуясь уравнениями


$\gamma^T L=\Lambda \gamma^T \Lambda^{-1}$ и


$\Lambda \gamma \Lambda^{-1}=L^{-1}\gamma$, множитель $\Lambda$


можно вынести, тогда все


слагаемые правой части, начиная со второго, в сумме дают нуль.


Это означает


$\nabla_\xi\psi'=\Lambda \nabla_\xi\psi.$





\section{Производная Ли спиноров в произвольных реперах}





Пусть дано некоторое поле репера $e(x)=(e_a{}^\alpha(x))$, к


которому привязаны контравариантное спинорное поле 1-го ранга


$\psi(x)$ и $\eta$ --- симметрическое преобразование $P$.


Бесконечно малое преобразование $x'=x+t\xi$ точку $x-t\xi$


переводит в точку $x$, репер $e(x-t\xi)$ --- в репер


$\widetilde e (x)=e(1+tK+\dotsb)$. При этом переходе считаем,


что $\psi (x-t\xi)$, $P(x-t\xi)$ и $g(x-t\xi)$ не меняются.


Они меняются при переходе от $\widetilde e(x)$ к $e(x)$. Образы


$\psi (x-t\xi)$, $P(x-t\xi)$ и $g(x-t\xi)$ обозначим


через $\widetilde \psi (x)$, $\widetilde P (x)$ и $\widetilde g(x)$.


Очевидно, $\widetilde g(x)=g(x)-t(K^T g+\partial_\xi g+gK)$.


Тогда в разложении $\widetilde P(x)=P(x)+t\delta P(x)$ вариация


$\delta P$ определяется из $\widetilde P \widetilde g \widetilde P=\eta$


как решение уравнения


$$


\delta P^TgP+P^Tg\delta P=P^T(K^Tg+gK+\partial_\xi g)P,


$$


которое эквивалентно уравнению (6) с


$\delta Q=P^2\eta^{-1}(K^Tg+gK+\partial_\xi g )P^2$.


Преобразование Лоренца $L$, соответствующее переходу от


$\widetilde e(x)$ к $e(x)$ вычисляется как


$$


L=\widetilde P^{-1}(1+tK)P(x-t\xi)=1-tP^{-1}(\delta P+


\partial_\xi P-KP).


$$


Тогда для производных Ли имеем


\begin{gather}


L_\xi g=\partial _\xi g+K^Tg+gK, \notag \\


L_\xi P=-\delta P, \\


L_\xi\psi=\partial_\xi \psi+\frac{1}{4}\gamma^TP^{-1}


(\delta P+\partial_\xi P-KP)\gamma\psi. \notag


\end{gather}


В ортрепере, когда $g=\eta$ и $P=1$, находим


\begin{gather*}


L_\xi\eta=2\eta K_C,\quad K_C=\frac{1}{2}(\eta^{-1}K^T\eta +K),\\


L_\xi 1=-K_C, \\


L_\xi\psi=\partial_\xi \psi-\frac{1}{4}\gamma^T K_A\gamma\psi.


\end{gather*}


Последнее выражение представляет производную Ли спинора (4),


найденную И.\,Косман. Соотношения (7) можно переписать в виде


\begin{gather*}


L_\xi g_{ab}=\nabla_b\xi _a+\nabla_a\xi_b, \\


L_\xi P=-\delta P, \\


L_\xi\psi=\nabla_\xi \psi+\frac{1}{4}\gamma^TP^{-1}(


\delta P -\nabla\xi P)\gamma\psi,


\end{gather*}


где $\nabla\xi=(\nabla_b\xi^a)$, а уравнение,


которому удовлетворяет $\delta P$, можно записать в виде


$$


P\delta P+\delta P P=P^2\eta^{-1}L_\xi gP^2.


$$





Чтобы утверждать, что $L_\xi \psi$ снова есть спинор, необходимо сначала


доказать, что $\frac{1}{4}\gamma^TP^{-1}\delta P\gamma \psi$


есть спинор, однако каково поведение этого выражения при произвольных


преобразованиях реперов, выяснить не удается. Удается однако объяснить,


почему для спинорных полей не выполняется свойство производных


Ли для векторных и тензорных полей, заключающееся в том, что


коммутатор производных Ли есть производная Ли вдоль коммутатора.





Если рассмотреть еще векторное поле $\mu^\alpha$ и обозначить


$M=-e^{-1}L_\mu e$, $N=-e^{-1}L_{[\xi,\mu]}e$, то можно показать,


что


\begin{gather}


[L_\xi,L_\mu]\psi=\partial_{[\xi,\psi]}\psi+\frac{1}{4}\gamma^T


(\partial_\mu K_A-\partial_\xi M_A+[K_A,M_A])\gamma\psi, \notag\\


L_{[\xi,\mu]}\psi=\Bigl\{[L_\xi,L_\mu]+\frac{1}{4}\gamma^T[K_C,M_C]


\gamma\Bigr\}\psi.


\end{gather}


Появление второго слагаемого в правой части последнего соотношения


можно объяснить следующим образом. Возьмем $P=1+tL_\mu 1=1-tM_C$


и по формуле (7) сосчитаем


$L_\xi (\psi+tL_\mu \psi)\vert_{P=1-tM_C}$.


Через $L_\xi\{L_\mu\psi\}$ будем обозначать значение


производной этого выражения по $t$ при $t=0.$ Если теперь


коммутатором двух производных Ли для $\psi$ считать


$L_\xi\{L_\mu\psi\}-L_\mu\{L_\xi\psi\},$ то для него мы в


точности получим правую часть (8).








\begin{thebibliography}{99}





\bibitem{1}


Картан Э. {\em Теория спиноров}. -- М.: ГИИЛ, -- 1947. -- 224 c.


\bibitem{2}


Van der Waerden B.L. {\em Spinoranalyse} // Nachr. Acad. Wiss. G\"ottingen:


Math.-phusik. kl. -- 1929. -- S.\,100--109.


\bibitem{3}


Fock V.A. {\em Geometrisierung der Diracshen Theorie des Elektrons} //


Zeitschz. f. Physik. -- 1929. -- Bd.\,57. -- S.\,261--277.


\bibitem{4}


Fock V.A., Ivanenko D.D. {\em Geometrie quantique lineaire et


deplacement parallele} // Compt. Rend. Acad. Sci. -- Paris,


1929. -- V.\,188. -- P.\,1470--1472.


\bibitem{5}


Infeld L., Van der Vaerden B.L. {\em Die Welengliechungen des Elektrons


in der allgemeinen Relativitatstheorie}


// Sitzber. Preuss. Acad. Wiss.


Physik.-math. k1. -- 1933. -- V.\,9. -- P.\,380--401.


\bibitem{6}


Rosenfeld L. {\em Sur le tenseur d'impulsion--energy} // Mem. Acad.


Roy. Belgique. -- 1940. -- V.\,18. -- fasc.\,6. -- P.\,1--30.


\bibitem{7}


Билялов Р.Ф. {\em Законы сохранения для спинорных полей на римановых


пространственно-временных многообразиях}


// Теор. и матем. физика. -- 1996. -- Т.\,108. -- \No\,2. -- C.\,306--314.


\bibitem{8}


Kosmann Y. {\em Derivees de Lie des spineurs} // Ann. math. pura ed


appl. -- 1972. -- T.\,91. -- \No\,4. -- P.\,317--395.


\bibitem{9}


Билялов Р.Ф. {\em Симметрический тензор энергии-импульса спинорных


полей} // Теор. и матем. физика. --


1992. -- T.\,90. -- \No\,3. -- C.\,369--379.


\bibitem{10}


Lichnerowicz A. {\em Champ de Dirac, champ de neitrino et


transformation C, P, T sur en espace--temps courbe}


// Ann. Inst. H.\,Poincare. -- 1964. -- V.\,1. -- \No\,3. -- P.\,233--290.


\bibitem{11}


Желнорович В.А. {\em Теория спиноров и ее применение в физике и


механике}. -- М.: Наука, 1982. -- 270 c.


\bibitem{12}


Кобаяси Ш., Номидзу К. {\em Основы дифференциальной геометрии}.


-- T.\,1. -- М.: Наука, 1981. -- 344 c.





\end{thebibliography}





\vskip 2cm





\post{\it Казанский государственный}{\it Поступила}


\post{\it университет}{10.12.1996}





\label{end}


\end{document}


